% !TEX root = ../ApproxDA-TransUNet.tex
\section{Experiments}
\label{sec:experiments}

\subsection{Datasets and Protocol}
\label{sec:protocol}

\textbf{Datasets.} \textit{Synapse}~\cite{landman2015miccai}: 30 abdominal CT scans covering 8 organs, with the standard 18/12 train/test split. \textit{Kvasir-SEG}~\cite{jha2020kvasir}: 1000 colonoscopy images with binary polyp masks, using the 880/120 train/test split recommended with the dataset. \textit{ISIC 2018}~\cite{codella2019isic}: dermoscopy images with binary lesion masks, using the official Task~1 training (2594), validation (100), and test (1000) sets.

\textbf{Model selection.} No checkpoint selection is performed: every model, including the re-trained DA-TransUNet, is evaluated at its final training epoch, as in TransUNet~\cite{chen2021transunet}. On ISIC 2018, $M$, $r$, $G$, and routing are selected on the official validation set. Synapse and Kvasir-SEG have no validation split; following common practice on Synapse~\cite{chen2021transunet,sun2024datransunet}, their configuration is taken from the ablations in Section~\ref{sec:understanding}, which report every configuration on the test set (see the limitations in Section~\ref{sec:limitations}). Each test set is otherwise used only for final evaluation.

\textbf{Training.} All models are trained for 300 epochs with ImageNet-21K pretrained R50-ViT-B/16 weights, $224{\times}224$ input, random rotation and flipping, batch size 24, polynomial learning-rate decay, and loss $\tfrac12\mathcal{L}_{\mathrm{CE}}+\tfrac12\mathcal{L}_{\mathrm{Dice}}$, on the same hardware (\tbd{one NVIDIA Tesla T4}). Following DA-TransUNet~\cite{sun2024datransunet}, Synapse uses SGD (learning rate 0.01, momentum 0.9, weight decay $10^{-4}$), and Kvasir-SEG and ISIC 2018 use Adam (learning rate $10^{-3}$, weight decay $10^{-4}$). We train longer than the original protocols (150 epochs on Synapse; 500 and 50 epochs on Kvasir-SEG and ISIC 2018) and use $224{\times}224$ rather than $256{\times}256$ input on all datasets, so that the same window sizes $M$ apply everywhere. DA-TransUNet is re-trained with exactly the same data, splits, and protocol. Unless stated otherwise, ApproxDA-TransUNet uses $r{=}32$ and PAM-only routing; $G$ ($=8$ by default) only matters for routings that use CAM.

\textbf{Metrics.} Synapse: mean Dice (DSC) and 95\% Hausdorff distance (HD95) over the 8 organs, computed with the TransUNet evaluation code. Kvasir-SEG and ISIC 2018: DSC and mean IoU (mIoU); we do not report HD95 for these datasets because their pixel spacing is not calibrated. \todo{main comparison: mean $\pm$ std over 3 seeds; Synapse per-organ DSC; paired test / bootstrap CI (experiments/statistical\_validation.py).}

\subsection{Main Results}
\label{sec:main}

\begin{table}[t]
  \caption{Segmentation Results. Rows marked $^\dagger$ are values reported in the literature under their original protocols and are shown for reference only; the matched comparison is between the last two rows.}
  \label{tab:main}
  \centering
  \footnotesize
  \setlength{\tabcolsep}{2.2pt}
  \begin{tabular}{@{}lcccccc@{}}
    \toprule
    & \multicolumn{2}{c}{Synapse} & \multicolumn{2}{c}{Kvasir-SEG} & \multicolumn{2}{c}{ISIC 2018} \\
    \cmidrule(lr){2-3}\cmidrule(lr){4-5}\cmidrule(lr){6-7}
    Method & DSC$\uparrow$ & HD95$\downarrow$ & DSC$\uparrow$ & mIoU$\uparrow$ & DSC$\uparrow$ & mIoU$\uparrow$ \\
    \midrule
    U-Net$^\dagger$~\cite{ronneberger2015unet}      & 76.85 & 39.70 & 88.22 & 80.12 & 87.22 & 81.14 \\
    U-Net++$^\dagger$~\cite{zhou2018unetpp}         & 76.91 & 36.93 & 86.57 & 77.67 & 87.30 & 81.33 \\
    Res-UNet$^\dagger$~\cite{diakogiannis2020resunet} & 76.95 & 38.44 & 77.85 & 66.04 & 83.32 & 76.51 \\
    MultiResUNet$^\dagger$~\cite{ibtehaz2020multiresunet} & 77.42 & 36.84 & -- & -- & -- & -- \\
    Att-UNet$^\dagger$~\cite{oktay2018attention}    & 77.77 & 36.02 & 86.61 & 78.01 & 87.60 & 81.51 \\
    \midrule
    TransUNet$^\dagger$~\cite{chen2021transunet}    & 77.48 & 31.69 & 87.91 & 80.03 & 88.78 & 82.63 \\
    UCTransNet$^\dagger$~\cite{wang2022uctransnet}  & 78.23 & 26.75 & --    & --    & --    & --    \\
    TransNorm$^\dagger$~\cite{azad2022transnorm}    & 78.40 & 30.25 & --    & --    & --    & --    \\
    MIM$^\dagger$~\cite{wang2022mim}                & 78.59 & 26.59 & --    & --    & --    & --    \\
    Swin-Unet$^\dagger$~\cite{cao2022swinunet}      & 79.13 & 21.55 & --    & --    & --    & --    \\
    DA-TransUNet$^\dagger$~\cite{sun2024datransunet} & 79.80 & 23.48 & 88.47 & --    & 88.88 & 82.78 \\
    \midrule
    DA-TransUNet (re-trained)                        & \tbd{--} & \tbd{--} & \tbd{88.44} & \tbd{81.70} & \tbd{88.43} & \tbd{80.68} \\
    ApproxDA-TransUNet (ours)                        & \tbd{80.94} & \tbd{27.49} & \tbd{90.17} & \tbd{84.27} & \tbd{89.55} & \tbd{--} \\
    \bottomrule
  \end{tabular}
  \\[2pt]
  \raggedright\scriptsize ApproxDA-TransUNet: PAM-only routing, $r{=}32$; $M$ = \tbd{Synapse 28, Kvasir-SEG 56, ISIC 2018 28}, selected from the ablations (Synapse, Kvasir-SEG) or on the validation set (ISIC 2018). Literature rows use different splits and training protocols (e.g., DA-TransUNet uses $256{\times}256$ input and Adam for 500/50 epochs on Kvasir-SEG/ISIC 2018).
\end{table}
% NOTE (remove before submission): current red numbers come from the single-seed
% BIBM runs (legacy block, test-selected checkpoints). The earlier DA-TransUNet
% Synapse re-run scored 72.03 DSC -- far below the reported 79.80; investigate
% before filling the re-trained row.

Table~\ref{tab:main} reports the results. Under matched training and evaluation, ApproxDA-TransUNet improves DSC over the re-trained DA-TransUNet on \tbd{all three} datasets (\tbd{$+1.73$} on Kvasir-SEG, \tbd{$+1.12$} on ISIC 2018, \todo{Synapse after re-run}). The gains are metric-specific rather than uniform: Synapse HD95 is \tbd{worse} than the DA-TransUNet and Swin-Unet values reported in the literature, and ISIC 2018 mIoU is \tbd{below} the reported DA-TransUNet value. Literature rows are included for context only and are not used to support our claims. Fig.~\ref{fig:qualitative} shows representative predictions.

\begin{figure}[t]
  \centering
  \includegraphics[width=\columnwidth]{figures/fig_qualitative_cross_task.png}
  \caption{Qualitative comparison (columns: input, ground truth, DA-TransUNet, ApproxDA-TransUNet). Top: Synapse ($M{=}\tbd{28}$); middle: Kvasir-SEG ($M{=}\tbd{56}$); bottom: ISIC 2018 ($M{=}\tbd{28}$).}
  \label{fig:qualitative}
\end{figure}
% TODO(figure): regenerate with the re-run (c16) checkpoints via
% experiments/ApproxDA-TransUNet/generate_qualitative_figure.py; the current image
% uses the BIBM legacy checkpoints (and its ISIC row used gate=learn, M=7).

\subsection{Computational Cost}
\label{sec:efficiency}

Table~\ref{tab:efficiency} compares parameters and FLOPs, counted with the same tool and input for all models. The unchanged backbone dominates end-to-end cost: the ViT encoder alone holds 85.1M parameters and 17.4 GFLOPs, and the CNN encoder and decoder most of the rest. The approximation reduces the cost of the attention blocks from \textbf{1.37} GFLOPs (exact attention) to \textbf{0.40} GFLOPs, below the 0.87 GFLOPs of the DA-TransUNet attention blocks even though ApproxDA-TransUNet uses one additional block at the $112^2$ skip. At the whole-network level the saving is small (29.73 vs.\ 30.20 GFLOPs), so we do not claim whole-network acceleration.

\begin{table}[t]
  \caption{Parameters and FLOPs (fvcore, one $1{\times}224{\times}224$ slice; one multiply--add counted as one FLOP). ``Attn.'' is the GFLOPs of the dual-attention blocks that are executed.}
  \label{tab:efficiency}
  \centering
  \footnotesize
  \setlength{\tabcolsep}{2.5pt}
  \begin{tabular}{@{}lccc@{}}
    \toprule
    Model & Params (M) & GFLOPs & Attn. \\
    \midrule
    DA-TransUNet~\cite{sun2024datransunet}$^\ast$ & 107.95 & 30.20 & 0.87 \\
    ApproxDA, exact ($M{=}112$, no proj.)         & 105.89 & 30.71 & 1.37 \\
    ApproxDA ($M{=}28$, $r{=}32$, PAM)            & 105.99 & 29.73 & 0.40 \\
    ApproxDA ($M{=}7$, $r{=}32$, learned)         & 106.51 & 30.08 & 0.74 \\
    \bottomrule
  \end{tabular}
  \\[2pt]
  \raggedright\scriptsize $^\ast$Public implementation; 1.40M of its parameters belong to attention blocks that are constructed but never executed. DA-TransUNet applies attention at three locations, ApproxDA-TransUNet at four.
\end{table}
% Numbers from experiments/audit_params_flops.py (static; independent of the re-runs).
